\chapter{Activity Diagrams}
\label{ch:activity}

\section{Purpose}

The activity diagrams model the work performed during a translation, from the
engineer's request through analysis to the point at which changes are written.

Where Chapter~\ref{ch:state} shows the states the tool rests in, these diagrams
show what happens between them. The two directions have their own diagrams
because they differ in shape, though both are organised around the same idea: a
queue of units to process, drained one at a time until empty.

\section{State Machine to Activity Diagram}
\label{sec:sa-activity}

This translation stays entirely within GENESYS: it reads a state machine from the
repository and produces an activity representing the same behaviour. No source
code is involved on either side, which is why neither of the code translations
covers it.

\umlfigure{state-to-activity}{State Machine to Activity activity diagram}{fig:sa-activity}

\subsection{Flow}

\begin{longtable}{L{0.30\textwidth}L{0.62\textwidth}}
	\caption{State to activity nodes}\label{tab:sa-nodes}\\
	\toprule
	\textbf{Node} & \textbf{Description} \\
	\midrule
	\endfirsthead
	\toprule
	\textbf{Node} & \textbf{Description} \\
	\midrule
	\endhead

	Connect to GENESYS & Establishes repository access, decomposed into the three
	steps that follow. The same sequence opens the export translation. \\
	\addlinespace
	Create Client, Log In, Open Project & The connection steps, drawn separately
	because each can fail independently. \\
	\addlinespace
	Select State Machine & The engineer chooses which machine to convert. \\
	\addlinespace
	Read States and Entry/Exit Behaviours & One half of \req{SA-1}. Retrieves the
	states together with the behaviours attached on entry and exit. \\
	\addlinespace
	Read Transitions, Triggers and Conditions & The other half of \req{SA-1}.
	Drawn in parallel with the states read, because neither depends on the
	other. \\
	\addlinespace
	Add States to Conversion Queue & Seeds the work queue. \\
	\addlinespace
	State Queue Empty? & The terminating condition. \\
	\addlinespace
	Take Next State from Queue & Removes the next state awaiting conversion. \\
	\addlinespace
	Create Action Node from Entry/Exit Behaviour & Produces the activity node
	corresponding to what the state does. \\
	\addlinespace
	Guarded Transitions? & Branches when a state's outgoing transitions carry
	guard conditions. \\
	\addlinespace
	Create Decision Node per Guard & Produces a decision node with one outgoing
	edge per guard, which is how a condition on a transition becomes a branch in an
	activity. \\
	\addlinespace
	Create Control Flow to Target States & Produces the edges joining this state's
	node to the nodes of the states it transitions to. \\
	\addlinespace
	Mark State as Reached & Records that the state appeared in the output, which
	is what makes \req{SA-3} answerable. \\
	\addlinespace
	Show All States and Transitions? & The option required by \req{SA-3}. \\
	\addlinespace
	Append Unreached States and Transitions & Adds the states and transitions that
	the traversal never arrived at, so the output covers the whole machine rather
	than only its reachable part. \\
	\addlinespace
	Review and Write & The shared sequence of
	Section~\ref{sec:shared-tail}. \\
	\bottomrule
\end{longtable}

\subsection{Design Decisions}

\paragraph{The two reads are parallel.}
States with their behaviours and transitions with their triggers and conditions
are independent retrievals, so they are drawn as a fork rather than a sequence.
Ordering them would assert a dependency the design does not have.

\paragraph{Guards become decision nodes, not separate states.}
A guard is a condition on a transition in a state machine, and a branch in an
activity. Creating one decision node per guarded state, rather than one per
transition, keeps the output at the abstraction a reader expects.

\paragraph{Reachability is tracked during the walk.}
\emph{Mark State as Reached} exists so that \emph{Append Unreached States and
Transitions} has something to work from. Without it, \req{SA-3} would require a
second traversal to discover what the first one missed.

\paragraph{The option is applied after the walk.}
\req{SA-3} is evaluated once the queue is empty rather than per state, so the
main loop is identical whichever way the option is set and only the final
appending step differs.

\section{Import Model from Code}
\label{sec:import-activity}

This direction reads source code and produces a GENESYS model. Work is organised
as a queue of methods, seeded with the entry point the engineer selects and
extended as the walk discovers calls.

\umlfigure{import-activity}{Import Model from Code activity diagram}{fig:import-activity}

\subsection{Flow}

\begin{longtable}{L{0.30\textwidth}L{0.62\textwidth}}
	\caption{Import activity nodes}\label{tab:import-nodes}\\
	\toprule
	\textbf{Node} & \textbf{Description} \\
	\midrule
	\endfirsthead
	\toprule
	\textbf{Node} & \textbf{Description} \\
	\midrule
	\endhead

	Select User Configurations & The engineer chooses the solution, the entry
	point and the depth to which calls are followed. \\
	\addlinespace
	Load Solution and Resolve Symbols & The solution is opened and symbols are
	resolved through the compiler's semantic model, so that a call is matched to
	the method it actually invokes rather than to a name. \\
	\addlinespace
	Symbols Resolved? & A quality gate. Branches when references cannot be
	resolved, typically from a project that does not build. \\
	\addlinespace
	Record Unresolved Symbols & Notes what could not be resolved and rejoins the
	main flow, so extraction continues with a partial but honestly labelled
	result. \\
	\addlinespace
	Seed Queue with Entry Point & Places the selected method into the work
	queue. \\
	\addlinespace
	Take Next Method from Queue & Removes the next method awaiting processing. The
	loop control point. \\
	\addlinespace
	Queue Empty? & The terminating condition. When no methods remain, control
	passes to the review and write sequence. \\
	\addlinespace
	Build Control Flow Graph & Produces the basic blocks and edges for one
	method. \\
	\addlinespace
	Reducible to Constructs? & A quality gate. Branches when control flow cannot
	be expressed as nested constructs, which early returns and breaks from nested
	loops can cause. \\
	\addlinespace
	Flag Method as Unstructurable & Records the method and the construct
	responsible, then returns to the queue, so one awkward method does not stop the
	run. \\
	\addlinespace
	Walk Control Flow Graph & Traverses the graph, dispatching on the structure of
	each branch and loop. \\
	\addlinespace
	Apply Abstraction Rules & Discards constructs below the significance
	threshold, which is what keeps the resulting diagram readable. \\
	\addlinespace
	Derive State Names from Code & Produces a name for each state from the
	identifiers and structure of the code it came from, as \req{CS-6} requires. It
	precedes entity creation because an entity cannot be created without its
	name. \\
	\addlinespace
	Map to GENESYS Entities & Produces the entities and relationships
	corresponding to what survived the abstraction. \\
	\addlinespace
	Enqueue Called Methods within Depth Limit & Extends the queue with methods the
	walk discovered, bounded by the depth the engineer selected, then returns to
	the queue. \\
	\addlinespace
	Report Unused Variables, Methods and Classes & Reports code the walk never
	reached, as \req{CS-5} requires. It sits after the queue drains because what is
	unused can only be known once everything reachable has been processed. \\
	\addlinespace
	Reconcile Renames and Moves & Matches code symbols to the model entities they
	produced on an earlier run rather than treating a rename as a deletion and a
	creation. Drawn dashed because \req{FS-3} is a stretch requirement. \\
	\addlinespace
	Review and Write & The shared sequence described in
	Section~\ref{sec:shared-tail}. \\
	\bottomrule
\end{longtable}

\subsection{Design Decisions}

\paragraph{A bounded queue rather than recursion.}
Following calls through a queue with an explicit depth limit keeps the traversal
finite, handles recursive and mutually recursive calls without special cases,
and means a method called from several places is processed once. That last
property is what satisfies \req{CS-4}, since recursion yields one queue item
rather than an expansion per occurrence.

\paragraph{Both quality gates continue rather than abort.}
Unresolved symbols and unstructurable methods are recorded and the run proceeds.
A partial result that states what it could not read is more useful than a failure,
and more honest than a result that silently omits it.

\paragraph{Abstraction after the walk, not during it.}
The walk produces a complete picture of the control flow, and the rules are
applied to that. Filtering during traversal would make the discarded material
unavailable for the counts reported at the end.

\section{Export Model to Code}
\label{sec:export-activity}

This direction reads a GENESYS model and produces source code. It opens with the
connection sequence, since the model lives in the repository rather than on
disk, then drains a queue of functions.

\umlfigure{export-activity}{Export Model to Code activity diagram}{fig:export-activity}

\subsection{Flow}

\begin{longtable}{L{0.30\textwidth}L{0.62\textwidth}}
	\caption{Export activity nodes}\label{tab:export-nodes}\\
	\toprule
	\textbf{Node} & \textbf{Description} \\
	\midrule
	\endfirsthead
	\toprule
	\textbf{Node} & \textbf{Description} \\
	\midrule
	\endhead

	Connect to GENESYS & Establishes access to the repository. Decomposed into the
	three steps that follow. \\
	\addlinespace
	Create Client & Instantiates the repository client, the only component
	referencing the vendor assemblies. \\
	\addlinespace
	Log In & Authenticates against the repository. \\
	\addlinespace
	Open Project & Selects the project whose model is to be translated. \\
	\addlinespace
	Get Functions & Retrieves the function entities from the opened project. \\
	\addlinespace
	Read Name, Number, Description & Reads each function's identifying attributes,
	including the description that becomes a documentation comment in the generated
	source. \\
	\addlinespace
	Add to Generation Queue & Places the functions into the work queue. \\
	\addlinespace
	Generate Code & The entry point to generation proper. \\
	\addlinespace
	Function Queue Empty & The terminating condition. \\
	\addlinespace
	Build Method & Emits the class for one function, with its function prototypes
	attached, as \req{AC-2} requires. Forks into three independent steps. \\
	\addlinespace
	Build Input/Output & Derives the prototype signatures from the function's
	inputs and outputs. \\
	\addlinespace
	Relate to Owning Class & Attaches the generated prototypes to the function
	class they belong to. \\
	\addlinespace
	Add Function Description & Emits the function's description as a documentation
	comment, carrying the intent recorded in the model into the source. \\
	\addlinespace
	Verify Generated Code Compiles & Checks the generated source against the
	target compiler, which is what makes \req{AC-4} a step in the flow rather than
	an unverified property of the output. \\
	\addlinespace
	Update in Repository & Writes the result. \\
	\bottomrule
\end{longtable}

\subsection{Design Decisions}

\paragraph{The connection sequence is explicit.}
Create client, log in and open project appear as distinct actions rather than
being folded into a single connect step, because each can fail independently and
because the client is the component that carries the vendor dependency.

\paragraph{Identifying attributes are read up front.}
Name, number and description are read when functions are retrieved rather than
when each method is built, so the queue holds everything generation needs and no
further repository round trip is required per function.

\paragraph{Method building forks rather than sequences.}
Inputs and outputs, owning class and description are drawn as parallel branches
because none depends on the others. Drawing them in sequence would assert an
ordering the design does not require.

\paragraph{Abstraction after the walk serves branch coverage.}
Applying the rules to a completed graph rather than during traversal is what
allows \req{CS-3} to be checked, since every branch is present before anything
is discarded.

\section{Shared Review and Write Sequence}
\label{sec:shared-tail}

Both directions end by presenting the proposed change set, waiting for the
engineer to accept or deny, writing accepted changes, and reporting the outcome.
This corresponds to the Awaiting Review, Writing and Showing Report states of
Figure~\ref{fig:state}.

Figure~\ref{fig:import-activity} represents this as a single \emph{Review and
Write} node, a call behaviour action invoking a separately defined activity,
rather than repeating the sequence in both diagrams.
